% Further NN-G036 refinement: root whole-proof PASS 2026-11-22T00:03:47Z.
\paragraph{The last mass-eight competitor.}
The remaining mass-eight base row has
\[
 (g_1,g_2,d_1,d_2,u_1,u_2)=(128,80,12,4,8,4).
\]
Thus $\delta_2=\psi$ and $U_2=D$. The integral divisibility
argument already used in NN-G035 gives
\[
 G_1=DWC,\qquad G_2=D(\psi+DC),\qquad
 C\equiv0\pmod2,\quad C(4)=4.
\]
For clarity, $DG_1=W(G_2-D\psi)$ and $\gcd(W,D)=1$
first give $G_1=WQ$ and $G_2=D(Q+\psi)$ integrally.
The other cell identity at zero forces $G_1(0)=0$,
so $Q(0)=0$ and $Q=DC$. Reducing $G_1=DWC$ modulo
two gives the displayed parity of $C$.
The irreducible mass-eight $W$ cannot divide $G_2$:
it is coprime to $D$, whereas the integral polynomial
$\psi+DC$ has value twenty at four, not divisible by eight.
Consequently $U_1G_2=W(G_1-D\delta_1)$ forces $W\mid U_1$.
The equal-mass lemma gives $U_1=W$. Cancelling now yields
\[
 \delta_1+\delta_2=(W-D)C.
\]
Write $C=2H$ with $H\in\mathbb Z[D]$ and $H(4)=2$.
The effective character
$F=(\delta_1+\delta_2)/2=(W-D)H$ has mass eight and
dominates $\psi$.
If $W=D+4$, then $F=4H$, so $H$ is effective;
domination forces its coefficients at both positive
exponents one and two to be at least one. Its mass
would therefore be at least four, contradicting two.
Otherwise $W=D+2T_j$ and $F=2T_jH$ has all Laurent
coefficients even. Its coefficients at exponents
$1,-1,2,-2$ must all be at least two. These exhaust
the mass eight, hence $F=2\psi$ and $T_jH=\psi$.
Since $T_j$ is monic of degree $j$, this requires
$j\le2$. But $T_1=D-2$ does not divide $\psi$, and
$T_2=D^2-4D+2$ is not the monic quadratic $\psi$.
This excludes the last mass-eight competitor.
Together with the preceding exclusions and NN-G034,
every competitive mass-eight closed table is the unique
$2744$ attainer. No bound on its initial support indices
was imposed.

\paragraph{The mass-ten block with equal action masses one hundred.}
The row is
\[
 (g_1,g_2,d_1,d_2,u_1,u_2)=(100,100,10,10,6,6).
\]
Thus both $U_i=D+2$. If $W$ is irreducible, its mass
ten divides neither $u_i$, so (G036.R2) forces
$W\mid G_1,G_2$. In the exceptional factorization
$W=(D-2)B$, the irreducible $B$ has mass five and
again divides neither $U_i$; it divides both $G_i$.
After cancelling $B$ in the cell identities, evaluation
at $D=2$, where $U_i(2)=4\ne0$, forces the remaining
factor $D-2$ into both $G_i$. Thus in every case
\[
 G_i=WP_i,\qquad P_i(4)=10,\qquad
 P_1\equiv0,\quad P_2\equiv D(D-1)\pmod2.
\]
The polynomial $P_2$ cannot be divisible by $D$,
since its quotient would have nonintegral value $10/4$
at four. The divided cell identities at zero then
force both $P_i(0)\ne0$ and give $W(0)^2=4$.
Consequently $W=D+2+2T_j$ with $j$ odd.
Write $\delta_i=\psi+2R_i$. Now $R_i$ is an effective
character of mass three, hence is either three or
$1+T_k$. A nonzero difference $E=R_1-R_2$ still
has leading coefficient $\pm1$ and $|E(0)|\le4$.
It is nonzero because the parity difference of
$P_1,P_2$ is nonzero.
Subtracting the divided identities gives exactly
\[
 \frac{D+2+T_j}{D}(P_1-P_2)=E.
\]
For $j=1$, the factor is two, contradicting the odd
leading coefficient of $E$. For odd $j\ge3$, it is
monic, with value $1+j^2$ at zero. Hence
$P_1-P_2=aD^2+bD+c$, where $a=\pm1$, $b$ is odd,
and $c$ is even: the degree follows from its leading
coefficient and its reduction $D(D-1)$.
The bound on $E(0)$ forces $c=0$.
Equal values of the two $P_i$ at four then give
$16a+4b=0$, contradicting that $b$ is odd.
The row is impossible.

\paragraph{The two mass-ten blocks with action masses eighty-four and one hundred.}
In either orientation the smaller action has
$\delta_-=\psi+2$, $U_-=D+2$, while the larger
has $\delta_+=\psi$ and $u_+=10$.
An irreducible $W$ of mass ten cannot divide $G_-$
of mass eighty-four; (G036.R2) consequently forces
$W\mid G_+$ and $W\mid U_+$, whence $U_+=W$
by the equal-mass lemma.
If $W=(D-2)B$, the mass-five irreducible factor
$B$ likewise cannot divide $G_-$. Thus $B\mid U_+$.
Writing $U_+=BR$, reduction gives $R\equiv D\pmod2$.
The leading coefficient argument in NN-G036 forces
$\deg R=1$; the mass-ten classification then gives
$U_+=W$ also in this case.
The larger-index identity now says $G_+=G_-+D\psi$.
Substituting in the smaller-index identity gives
\[
 (W-D-2)G_-=D\bigl(W(\psi+2)+(D+2)\psi\bigr).
\]
For $W=D+6$ this would make the leading coefficient
of $G_-$ equal to one half. Otherwise
$W=D+2+2T_j$ and
\[
 G_-=D(\psi+2)+\frac{N(D)}{T_j(D)},\qquad
 N=D(D+2)(D^2-3D+1).
\]
The same integral divisibility obstruction as in
(G036.R6) applies: $T_j\mid N$ requires $j\le4$,
but $N(2)=-8$, $T_1\mid T_3$,
$N\bmod T_2=16D-10$, and $N\ne T_4$.
Both orientations are excluded, with every auxiliary
increment allowed by their scalar relaxation.
\hfill$\square$

\paragraph{Scope and geometric content.}
These four base-row exclusions concern the ordinary
dimensions of actual graded $C_2$-representation
structure objects in the matrix-kernel category of
NN-G010, with the same symmetric closed multiplication
table and character constraints. In particular, within
the mass-eight gluing stratum the actual smallest
effective lift is the $2744$-pair object of NN-G033,
uniquely within that presentation. The integral torus
character interpretation is SGA~3, I, 4.7.2--4.7.3;
kernel composition is the supplied archive transcription,
cote~2, scan~17. No new claim about a canonical GCT
representation family or a different multiplication
table follows.
